\documentclass[10pt,a4paper]{amsart}
\usepackage[margin=19mm]{geometry}
\usepackage{amsmath,amssymb,mathtools}
\usepackage[scr=rsfs,cal=euler]{mathalfa}
\usepackage{xfrac}
\usepackage[unicode,hidelinks,bookmarks=false]{hyperref}
\newcommand{\ie}{i.e.\ }
\newcommand{\cf}{cf.\ }
\newcommand{\resp}{respectively\ }
\newcommand{\Subsection}{\subsection}
\providecommand{\nobreakdash}{\nobreak\hbox{-}}
\setlength{\emergencystretch}{2em}
\multlinegap0pt
\usepackage{amssymb}
\usepackage{enumerate}
\usepackage[shortlabels]{enumitem}
\usepackage{cancel}
\usepackage{tikz}
\usepackage{booktabs}
\usepackage{nicefrac}
\usepackage{multirow}
\usepackage{xfrac}
\usepackage{caption}
\usepackage{bbm}
\usepackage{float}
\newtheorem{lemma}{Lemma}
\title{A Practical Adaptive Subgame Perfect Gradient Method}
\author{Alan Luner, Benjamin Grimmer}
\date{Source: arXiv:2510.21617v2 (CC BY 4.0)}
\begin{document}
\maketitle

\noindent\emph{Source and licence.} A Practical Adaptive Subgame Perfect Gradient Method. Version arXiv:2510.21617v2 (CC BY 4.0), distributed by arXiv under the Creative Commons Attribution 4.0 International licence, http://creativecommons.org/licenses/by/4.0/, as recorded in the arXiv licence metadata for this exact version and shown on the abstract page https://arxiv.org/abs/2510.21617v2. Full licence notice: \texttt{licenses/arxiv-2510.21617-CC-BY-4.0.txt}.\par\smallskip

\noindent\textit{Editorial scope.} Counted unit: the article's lemma Lem:BetaInequalities (source0.tex lines 1065--1071). The article's complete original proof of it is delivered with it (source0.tex lines 1073--1104, one \texttt{proof} environment of 32 lines). No mathematics is written, repaired or re-derived: the statement and the proof are the article's own lines, in the article's own notation, and no retained line is substituted or re-flowed.  No further source-local context is needed: every cross-reference of every retained line resolves inside the two delivered spans.  Nothing was omitted from any retained span, so the package carries no omission record.  No retained line contains a citation, so the wrapper carries no bibliography and no bibliography entry.  No retained line contains a figure environment, an image-inclusion command or drawing code, so no image is delivered and the package declares no asset dependency.  Editorial formatting only, in addition: an AMS \texttt{amsart} preamble that restates the article's own declarations and macros used by the retained lines, namely the environments \texttt{lemma} on separate counters, together with the standard packages those lines need.  The retained environments print in this document as Lemma 1 in order of appearance, whereas the article prints the same items with the numbering of its own full document.  The complete native source of the article as distributed in the arXiv e-print for this exact version is bundled unchanged as \texttt{sources/187737/source0.tex}.
\begin{lemma} \label{Lem:BetaInequalities}
    Let $j > i \geq n$. Then the following hold
    \begin{align}
        \beta_i &\leq \frac{1}{\phi_n (2\psi_i + 1)} \label{Eqn:BetaLemma1} \\
        (2\psi_j + 1) \beta_j & \leq (2 \psi_i - 1) \beta_i. \label{Eqn:BetaLemma2}
    \end{align}
\end{lemma}

\begin{proof}
    Note from our definition of $\psi_i$ that $\psi_i^2 = 2 \tau_{n,i} - \psi_i = \phi_i + \tau_{n,i}$ for $n \leq i < N$ and $\psi_N^2 = \tau_{n,N} - \psi_N = \phi_N$. From this we have
    \begin{equation*}
        \tau_{n,i}(2\psi_i - 1) = \begin{cases}
            \phi_i(2\psi_i+1)+\psi_i^2 \quad & \text{if } i \in [n, N-1] \\
            \phi_N(2\psi_N+1)+\psi_N^2 - \tau_{n,N} \quad & \text{if } i = N.
        \end{cases}
    \end{equation*}
Using this relation, we can write
\begin{align*}
    \phi_{i+1}(2 \psi_{i+1} + 1) \beta_{i+1} &= 1 + \sum_{\ell = n}^i \psi_\ell^2 \beta_\ell \\
    & = 1 + \sum_{\ell = n}^{i-1} \psi_\ell^2 \beta_\ell + \psi_i^2 \beta_i \\
    & = \phi_{i}(2 \psi_{i} + 1) \beta_{i} + \psi_i^2 \beta_i \\
    & = \tau_{n,i}(2 \psi_i - 1) \beta_i.
\end{align*}
Recognizing that $\phi_{i+1} = \tau_{n,i}$, we rearrange to obtain the identity 
\begin{equation} \label{Eqn:BetaIdentity}
    (2 \psi_{i+1} + 1) \beta_{i+1} = (2 \psi_i -1) \beta_i.
\end{equation}

We now prove \eqref{Eqn:BetaLemma1} by induction. By definition, $\beta_n = \frac{1}{\phi_n(2 \psi_n+1)}$, so our base case holds. Then for $i>n$, using \eqref{Eqn:BetaIdentity} we can write
\begin{align*}
    \beta_i = \frac{(2 \psi_{i-1} -1)\beta_{i-1}}{2 \psi_i + 1} \leq \frac{2 \psi_{i-1} - 1}{2\psi_i + 1} \frac{1}{\phi_n (2 \psi_{i-1} + 1)} \leq \frac{1}{\phi_n (2 \psi_{i-1} + 1)}
\end{align*}
where we use our induction hypothesis and then the fact that $\psi_i$ is monotonically increasing.

For the second inequality \eqref{Eqn:BetaLemma2}, applying our identity \eqref{Eqn:BetaIdentity} yields
\begin{equation*}
    (2 \psi_j + 1)\beta_j = (2 \psi_{j-1} - 1) \beta_{j-1} \leq (2 \psi_{j-1} + 1) \beta_{j-1}
\end{equation*}
and the result follows from induction.
\end{proof}

\end{document}
